\section{Conclusiones}
\noindent Los microcuásares y los cuásares manifiestan representan, a diferentes escalas, el mismo tipo de fenómenos físicos. Dado que su radio gravitacional y escala temporal dinámica escalan linealmente con la masa del agujero negro, los microcuásares sufren los mismos procesos que a los cuásares les tomarían siglos, e incluso miles de años, en escalas humanas.  \\

\noindent El estudio cinemático y evolutivo de los microcuásares permite reconstruir cómo mueren las estrellas masivas, demostrando que la masa del remanente determina si el colapso ocurre de forma explosiva o silenciosa. Tras su fase de alta energía, estos sistemas culminan su ciclo como binarias compactas destinadas a fusionarse mediante la emisión de ondas gravitacionales. \\

\noindent Los estados espectrales de los microcuásares reflejan distintas geometrías de acreción: un disco delgado que llega hasta el $R_{\text{ISCO}}$ en el estado blando, frente a un estado duro dominado por Comptonización térmica. Para explicar este último se plantean dos modelos principales: un disco frío truncado con un flujo interior caliente (ADAF/RIAF), o una corona dinámica eyectada (\textit{outflowing} ADC) vinculada a la base del \textit{jet}. \\

\noindent Los estallidos de las binarias de rayos X describen un ciclo de histéresis en forma de <<Q>> gobernado por la asimetría entre la evaporación del disco (estado duro) y el enfriamiento Compton (estado blando). Esto impone una alternancia excluyente en la eyección de materia: el estado duro emite jets polares relativistas (divididos en dos pistas según su eficiencia radiativa), mientras que el estado blando los suprime y activa vientos ecuatoriales ionizados. Al transicionar entre ambos estados, el colapso del disco hasta la última órbita estable ($R_{\text{ISCO}}$) apaga el \textit{jet} continuo y lanza eyecciones discretas ultrarrelativistas. 